% Figure: power curves of a one-sample z-test as a function of the
% standardised effect size, for three sample sizes. Coordinates are
% precomputed values of Phi(d*sqrt(n) - z_{0.975}) at alpha=0.05.
\begin{figure}[htbp]
\centering
\begin{tikzpicture}
\begin{axis}[
    width=11cm, height=6cm,
    xmin=0, xmax=1.0,
    ymin=0, ymax=1.02,
    axis lines=left,
    xlabel={Standardised effect size $d = (\mu_1 - \mu_0)/\sigma$},
    ylabel={Power $1 - \beta$},
    every axis plot/.append style={very thick, mark=none},
    legend style={font=\small, draw=none, fill=none, at={(0.98, 0.05)}, anchor=south east},
    xtick={0, 0.2, 0.4, 0.5, 0.6, 0.8, 1.0},
    ytick={0, 0.2, 0.5, 0.8, 0.95, 1.0},
    grid=both,
    grid style={dashed, draw=enggray!20},
]
% n = 20: power at d in {0.0, 0.1, ..., 1.0}
\addplot[engblue, smooth] coordinates {
  (0.00, 0.025) (0.10, 0.060) (0.20, 0.130) (0.30, 0.244)
  (0.40, 0.396) (0.50, 0.564) (0.60, 0.718) (0.70, 0.838)
  (0.80, 0.918) (0.90, 0.963) (1.00, 0.985)
};
\addlegendentry{$n = 20$}

% n = 50: power at d in {0.0, 0.1, ..., 1.0}
\addplot[engred, smooth] coordinates {
  (0.00, 0.025) (0.10, 0.097) (0.20, 0.293) (0.30, 0.567)
  (0.40, 0.807) (0.50, 0.940) (0.60, 0.986) (0.70, 0.998)
  (0.80, 1.000) (0.90, 1.000) (1.00, 1.000)
};
\addlegendentry{$n = 50$}

% n = 100: power at d in {0.0, 0.1, ..., 1.0}
\addplot[englime!70!black, dashed, smooth] coordinates {
  (0.00, 0.025) (0.10, 0.170) (0.20, 0.516) (0.30, 0.851)
  (0.40, 0.979) (0.50, 0.999) (0.60, 1.000) (0.70, 1.000)
  (0.80, 1.000) (0.90, 1.000) (1.00, 1.000)
};
\addlegendentry{$n = 100$}

\addplot[engblack, dotted, mark=none] coordinates {(0, 0.8) (1.0, 0.8)};
\end{axis}
\end{tikzpicture}
\caption[Power curves of the one-sample $z$-test]{Power of the
one-sample $z$-test at significance $\alpha = 0.05$ as a function of
the standardised effect size $d = (\mu_1 - \mu_0)/\sigma$. The
dotted horizontal line marks the conventional $80\,\%$-power
benchmark. The required effect size for $80\,\%$ power decreases
roughly as $1/\sqrt{n}$. Power is the design quantity the
engineer reports alongside $\alpha$; a study with high power against
a small effect is the converse design of a study with low power
against a large one.}
\label{fig:vol02:ch14:power-curve}
\end{figure}
